\subsubsection{Applications to Elliptic PDEs}\label{subsec:implicit-function-elliptic-applications}

We want to apply the results of the previous subsection to nonlinear elliptic PDEs. In order to do this, we have to set up such problems as abstract nonlinear equations in Banach space. The following lemma is crucial for this.

\begin{Lemma}
\label{lem:sobolev-pointwise-composition-differentiability}
Let $\Omega$ be a bounded domain in $\mathbb{R}^n$ with smooth boundary. Moreover, let $f:\mathbb{R}^k\to\mathbb{R}$ be of class $C^{m+1}$, where $m>n/2$. Then $f$, interpreted pointwise, induces a continuously differentiable mapping from $(H^m(\Omega))^k$ into $H^m(\Omega)$.
\end{Lemma}

\begingroup
\renewcommand{\endProof}{\par\nobreak\hfill$\blacksquare$\par\vspace{0.1cm}}
\begin{Proof}
By the Sobolev embedding theorem, we have $H^m(\Omega)\subset C(\overline\Omega)$. Hence we have $f(\bm{u})\in C(\overline\Omega)\subset L^2(\Omega)$ for every $\bm{u}\in(H^m(\Omega))^k$. Moreover, if $\bm{u}$ is in $(C^m(\overline\Omega))^k$, we can obtain the derivatives of $f(\bm{u})$ by the chain rule, and in the general case, we can use approximation by smooth functions. For example, let $\bm{u}=\lim_{n\to\infty}\bm{u}_n$ in the topology of $(H^m(\Omega))^k$, where the $\bm{u}_n$ are smooth. We find
\begin{equation}\label{eq:sobolev-composition-distributional-chain-rule}
\begin{aligned}
\left(f(\bm{u}),\frac{\partial\phi}{\partial x_i}\right)
&=\lim_{n\to\infty}\left(f(\bm{u}_n),\frac{\partial\phi}{\partial x_i}\right)\\
&=-\lim_{n\to\infty}\left(\nabla f(\bm{u}_n)\frac{\partial\bm{u}_n}{\partial x_i},\phi\right)\\
&=-\left(\nabla f(\bm{u})\frac{\partial\bm{u}}{\partial x_i},\phi\right),
\end{aligned}
\end{equation}
i.e., the chain rule applies for differentiating $f(\bm{u})$. In a similar fashion, we can deal with higher derivatives. Note that all derivatives of $f(\bm{u})$ have the form of a product involving a derivative of $f$ and derivatives of $\bm{u}$. The first factor is in $C(\overline\Omega)$, while any $l$th derivative of $\bm{u}$ lies in $H^{m-l}(\Omega)$, which embeds into $L^{2n/(n-2(m-l))}(\Omega)$ if $m-l<n/2$. We can use this fact and Hölder\textquotesingle s inequality to show that all derivatives of $f(\bm{u})$ up to order $m$ are in $L^2(\Omega)$; moreover, it is clear from this argument that $f$ is actually continuous from $(H^m(\Omega))^k$ into $H^m(\Omega)$. The differentiability of $f$ follows along similar lines by exploiting the relation
\begin{equation}\label{eq:sobolev-composition-integrated-difference-formula}
f(\bm{u})-f(\bm{v})=\int_0^1\nabla f(\bm{v}+\theta(\bm{u}-\bm{v}))\cdot(\bm{u}-\bm{v})\,\mathrm{d}\theta.
\end{equation}
We leave it to the reader to fill in some more details of the proof (Problem~\ref{ex:sobolev-composition-lemma-details}).
\end{Proof}
\endgroup
